\documentclass[11pt,a4paper]{article}

\usepackage[margin=1in]{geometry}
\usepackage{amsmath}
\usepackage{enumitem}
\usepackage{hyperref}
\usepackage{listings}
\usepackage{xcolor}

\title{Disease Symptom Prediction using Retrieval-Augmented Generation (RAG)}
\author{AIO Vietnam - Module 1 Project Proposal}
\date{\today}

\begin{document}

\maketitle

\section{Project Goal}

Build a lightweight and low-cost RAG system that accepts a user's symptom query and returns:

\begin{itemize}
    \item Possible diseases
    \item Disease descriptions
    \item Severity information
    \item Recommended precautions
\end{itemize}

The objective is educational and demonstrational rather than clinical diagnosis.

\section{System Architecture}

The proposed MVP architecture is shown below:

\begin{verbatim}
+------------------+
| User Query       |
| "I have fever,   |
| cough, fatigue"  |
+--------+---------+
|
v
+-------------------+
| Query Embedding   |
| BGE-small-en-v1.5 |  (in app)
+--------+----------+
         |
         v
+-----------------------------------+
| OpenSearch (Aiven, managed)       |
|   hybrid query + search pipeline  |
|  +-----------+   +-------------+  |
|  | BM25      |   | k-NN vector |  |
|  | (text)    |   | (384-dim)   |  |
|  +-----+-----+   +------+------+  |
|        |                |        |
|        +-------+--------+        |
|                v                 |
|    RRF (score-ranker) -> top 20  |
+----------------+------------------+
                 |
                 v
+------------------+
| BGE Reranker     |
| top 20 -> top 5  |  (in app)
+--------+---------+
         |
         v
+------------------+
| LLM              |
| Qwen3 8B         |
+--------+---------+
         |
         v
+------------------+
| Final Response   |
+------------------+
\end{verbatim}

\section{Data Preparation}

Each disease is stored as a single document. The source data follows the
classic disease--symptom dataset format:

\begin{lstlisting}[basicstyle=\ttfamily\small]
{
"disease": "Influenza",
"symptoms": ["fever", "cough", "fatigue"],
"severity": "medium",
"description": "...",
"precautions": "..."
}
\end{lstlisting}

\subsection{Three Logical Views of a Document}

The key insight for \emph{symptom-to-disease} retrieval is that the user query
is a list of symptoms, so the searchable text must be symptom-dense. The
description and precautions are mostly \emph{output} fields, not \emph{search}
fields. Each record is therefore projected into three distinct views:

\begin{enumerate}
    \item \textbf{Embedded text} (used to compute the dense vector stored in the
          \texttt{embedding} \texttt{knn\_vector} field): symptom-first natural
          language.
    \item \textbf{Keyword text} (indexed as an analyzed \texttt{text} field for
          BM25): raw symptom tokens, the disease name, and synonyms.
    \item \textbf{Metadata payload} (returned in \texttt{\_source}, not searched):
          severity, description, precautions, and provenance.
\end{enumerate}

All three views live in a single OpenSearch document, so there is one source of
truth and no cross-store synchronization.

\subsection{Embedded Text Template}

A symptom-first template keeps the embedding centroid close to symptom queries:

\begin{lstlisting}[basicstyle=\ttfamily\small]
def build_embed_text(rec: dict) -> str:
    symptoms = ", ".join(rec["symptoms"])
    return (
        f"Disease: {rec['disease']}. "
        f"Symptoms: {symptoms}. "
        f"{rec['description']}"
    )
\end{lstlisting}

\subsection{Symptom Normalization}

The raw dataset uses underscores (\texttt{skin\_rash}, \texttt{high\_fever}).
These are converted to natural language (\texttt{skin rash},
\texttt{high fever}) before embedding, because the embedding model handles
natural phrases far better than \texttt{snake\_case} tokens.

\subsection{Query-Time Synonym Expansion}

A small synonym map is applied to the user query (e.g.\ ``tired''
$\rightarrow$ ``fatigue'', ``throwing up'' $\rightarrow$ ``vomiting''). At this
dataset scale, this improves recall more than swapping the embedding model.

\subsection{Document Fields (OpenSearch Mapping)}

Everything needed to render the final answer is stored in the document
\texttt{\_source}, so source data is never re-parsed after retrieval. Severity
and precautions are deliberately kept \emph{out} of the embedded text to avoid
adding noise to similarity; severity is more useful as a filter
(\texttt{keyword} type) or a post-retrieval field.

\begin{center}
    \begin{tabular}{|l|l|l|}
        \hline
        Field                  & OpenSearch type   & Purpose                     \\
        \hline
        \texttt{doc\_id}       & keyword           & stable id / provenance join \\
        \texttt{disease}       & text + keyword    & display and grouping        \\
        \texttt{symptoms}      & keyword[]         & reranking and explanation   \\
        \texttt{severity}      & keyword           & display and filtering       \\
        \texttt{description}   & text              & LLM context                 \\
        \texttt{precautions}   & keyword[]         & LLM context                 \\
        \texttt{keyword\_text} & text              & BM25 lexical field          \\
        \texttt{embedding}     & knn\_vector (384) & dense vector (cosine)       \\
        \texttt{source}        & keyword           & provenance / citation       \\
        \hline
    \end{tabular}
\end{center}

\section{Embedding Model}

\subsection{Recommendation}

For Module 1, use:

\begin{itemize}
    \item BGE-small-en-v1.5
\end{itemize}

Advantages:

\begin{itemize}
    \item Fast inference
    \item Small memory footprint
    \item Easy deployment
    \item Good retrieval quality
\end{itemize}

\subsection{Embedding Strategy}

Three practices materially improve retrieval quality with BGE:

\begin{itemize}
    \item \textbf{Asymmetric query prefix.} BGE is trained with an instruction
          prefix. Queries are embedded as ``Represent this sentence for
          searching relevant passages: '' $+$ query, while documents are
          embedded \emph{without} the prefix.
    \item \textbf{Normalized vectors + cosine similarity.} Embeddings are
          L2-normalized and the \texttt{knn\_vector} field uses the
          \texttt{cosinesimil} space.
    \item \textbf{Exact / approximate search at MVP scale.} The corpus is on the
          order of tens of diseases, so exact k-NN is more than adequate;
          OpenSearch's HNSW approximate index only matters once the corpus grows
          much larger.
\end{itemize}

The embedding is computed in the application (the embedding model is not hosted
inside OpenSearch); the resulting 384-dimensional vector is sent into the
\texttt{knn} sub-query at search time and stored in the \texttt{embedding}
field at index time.

\begin{lstlisting}[basicstyle=\ttfamily\small]
from sentence_transformers import SentenceTransformer
model = SentenceTransformer("BAAI/bge-small-en-v1.5")

doc_emb = model.encode(doc_texts, normalize_embeddings=True)
query_emb = model.encode(
    [query], normalize_embeddings=True,
    prompt="Represent this sentence for searching relevant passages: ",
)
\end{lstlisting}

\subsection{Future Upgrade}

Possible future embeddings:

\begin{itemize}
    \item BioBERT
    \item PubMedBERT
\end{itemize}

\section{Search Engine}

A single managed \textbf{OpenSearch} service handles both retrieval modes: BM25
lexical search on analyzed \texttt{text} fields and dense vector search via the
bundled k-NN plugin. This removes the need for a separate vector database --
one engine, one source of truth, and no cross-store synchronization.

Advantages:

\begin{itemize}
    \item BM25 lexical search and k-NN dense vector search in one engine
    \item Server-side hybrid fusion (including RRF) via a search pipeline, so no
          application-level merge code is required
    \item Metadata filtering (e.g.\ \texttt{severity}) in the same query
    \item Available on the Aiven free tier (4\,GB RAM, 20\,GB storage), so the
          MVP is fully managed at no cost
\end{itemize}

\subsection{Managed Deployment (Aiven)}

OpenSearch runs as a managed Aiven service from the start, so there is no
container or volume to operate and the seeded index persists across runs.
Connection credentials are provided to the application via environment
variables / secrets and are never committed to source control. The free tier is
a single node (no high availability), which is acceptable for the MVP; a paid
plan adds replicas for production.

\subsection{Index Mapping}

The index defines the lexical field, the k-NN vector field, and the metadata
fields returned in \texttt{\_source}:

\begin{lstlisting}[basicstyle=\ttfamily\small]
PUT /diseases
{
  "settings": { "index": { "knn": true } },
  "mappings": {
    "properties": {
      "doc_id":       { "type": "keyword" },
      "disease":      { "type": "text",
                        "fields": { "raw": { "type": "keyword" } } },
      "symptoms":     { "type": "keyword" },
      "severity":     { "type": "keyword" },
      "description":  { "type": "text" },
      "precautions":  { "type": "keyword" },
      "keyword_text": { "type": "text" },
      "embedding": {
        "type": "knn_vector",
        "dimension": 384,
        "space_type": "cosinesimil"
      }
    }
  }
}
\end{lstlisting}

\subsection{Idempotent Ingestion}

Documents are indexed with a deterministic id (\texttt{\_id = doc\_id}), so
re-running ingestion upserts rather than duplicates -- the load is idempotent
and safe to repeat:

\begin{lstlisting}[basicstyle=\ttfamily\small]
POST /diseases/_doc/influenza        # _id = doc_id => upsert, no duplicates
{
  "doc_id": "influenza",
  "disease": "Influenza",
  "symptoms": ["fever", "cough", "fatigue"],
  "severity": "medium",
  "description": "...",
  "precautions": ["rest", "fluids"],
  "keyword_text": "influenza fever cough fatigue ...",
  "embedding": [0.013, -0.052, ...]   # 384-dim BGE vector (app-computed)
}
\end{lstlisting}

\section{Retrieval Strategy}

A hybrid retrieval approach is recommended.

\subsection{BM25 Retrieval}

A standard \texttt{match} query on the analyzed \texttt{keyword\_text} field,
scored with OpenSearch's default Okapi BM25.

Advantages:

\begin{itemize}
    \item Exact keyword matching
    \item Effective for symptom names
\end{itemize}

\subsection{Vector Retrieval (k-NN)}

A \texttt{knn} query on the \texttt{embedding} field using cosine similarity.

Advantages:

\begin{itemize}
    \item Semantic search
    \item Handles paraphrases and synonyms
\end{itemize}

\subsection{Hybrid Search Pipeline}

Both retrievers run inside OpenSearch within a single \texttt{hybrid} query, and
their result sets are fused server-side by a search pipeline using
\textbf{Reciprocal Rank Fusion (RRF)} (the \texttt{score-ranker-processor}). RRF
is preferred over weighted score mixing because BM25 scores are unbounded and
corpus-dependent while cosine similarity is bounded in $[-1, 1]$; RRF uses only
ranks, so no per-dataset weight calibration is required.

\[
    \mathrm{score}(d) = \sum_{r \in \{\text{bm25},\,\text{knn}\}}
    \frac{1}{k + \mathrm{rank}_r(d)}, \quad k = 60
\]

\begin{verbatim}
query
 |-- BM25 (match)        -> ranked list --|
 |-- k-NN (cosine)       -> ranked list --|
                                          v
                          RRF (score-ranker processor)
                                          v
                                     top 20 docs
\end{verbatim}

Define the fusion pipeline once, then attach it at search time:

\begin{lstlisting}[basicstyle=\ttfamily\small]
PUT /_search/pipeline/hybrid-rrf
{
  "phase_results_processors": [
    { "score-ranker-processor": { "combination": { "technique": "rrf" } } }
  ]
}

GET /diseases/_search?search_pipeline=hybrid-rrf
{
  "size": 20,
  "query": {
    "hybrid": {
      "queries": [
        { "match": { "keyword_text": "fever cough fatigue" } },
        { "knn": { "embedding": { "vector": [/* 384-dim */], "k": 20 } } }
      ]
    }
  }
}
\end{lstlisting}

The application then reranks these 20 candidates with \texttt{bge-reranker-base}
and passes the top 5 to the LLM.

\section{Reranking}

Use:

\begin{itemize}
    \item bge-reranker-base
\end{itemize}

Purpose:

\begin{itemize}
    \item Improve relevance ranking
    \item Filter noisy retrieval results
    \item Select top disease candidates
\end{itemize}

Pipeline:

\begin{verbatim}
Top 20 Documents
|
v
BGE Reranker
|
v
Top 5 Documents
\end{verbatim}

\section{Large Language Model}

\subsection{Recommended Model}

Qwen3 8B

Reasons:

\begin{itemize}
    \item Open-source
    \item Low deployment cost
    \item Good reasoning capability
    \item Suitable for local inference
\end{itemize}

\subsection{LLM Responsibilities}

The LLM should:

\begin{itemize}
    \item Summarize retrieved information
    \item Explain disease similarities
    \item Present possible diseases
    \item Recommend precautions
    \item Express uncertainty when appropriate
\end{itemize}

The LLM should not be treated as a medical diagnostic system.

\section{Technology Stack}

\begin{center}
    \begin{tabular}{|l|l|}
        \hline
        Component      & Technology                         \\
        \hline
        Embedding      & BGE-small-en-v1.5 (384-dim)        \\
        Search Engine  & OpenSearch (Aiven, managed)        \\
        Vector Search  & OpenSearch k-NN (cosine)           \\
        Keyword Search & OpenSearch BM25                    \\
        Fusion         & RRF (score-ranker search pipeline) \\
        Reranker       & bge-reranker-base                  \\
        LLM            & Qwen3 8B                           \\
        Framework      & LangChain / LlamaIndex (optional)  \\
        \hline
    \end{tabular}
\end{center}

\section{Scalability Roadmap}

\subsection{Phase 1 (MVP)}

\begin{itemize}
    \item OpenSearch on the Aiven free tier (BM25 + k-NN)
    \item Hybrid retrieval with server-side RRF fusion
    \item BGE-small-en-v1.5 embeddings
    \item bge-reranker-base
    \item Qwen3 8B
\end{itemize}

\subsection{Phase 2 (Production)}

Because OpenSearch is used from the start, scaling is mostly a matter of moving
to a larger plan and tuning the same engine rather than migrating systems:

\begin{itemize}
    \item Paid OpenSearch plan with replicas (high availability) and HNSW tuning
          for larger corpora
    \item BioBERT or PubMedBERT embeddings
    \item Larger reranker models
    \item GPT-5 or larger open-source LLMs
    \item Monitoring and feedback loops
\end{itemize}

\section{Conclusion}

For Module 1, the following architecture provides the best balance between simplicity, cost, and scalability:

\begin{verbatim}
User Query
|
v
Embedding (BGE-small, in app)
|
v
OpenSearch hybrid: BM25 + k-NN  ->  RRF (top 20)
|
v
Reranker (bge-reranker-base, top 5)
|
v
Qwen3 8B
|
v
Response
\end{verbatim}

This architecture is easy to implement, inexpensive to operate, and can be scaled incrementally as project requirements grow.

\end{document}
